\documentclass[12pt]{article}

\usepackage[spanish]{babel}
\usepackage{amsmath}
\usepackage{amssymb}

\title{Cálculo de la cantidad óptima de hielo para enfriar un volúmen de agua}
\author{Jorge Alberto Suárez Saldaña}
\date{\today}

\begin{document}
\maketitle
\section{Introducción}
Tenemos un problema interesante: En Marte Boba Shop tenemos un dispensador de agua con capacidad de 10 litros que está
pensado para mantener dicha agua fría.
En una ocación, mientras rellenaba dicho dispensador con hielos, me pregunté: 
Si le agrego un cubo de hielo al dispensador con agua, el hielo se derretirá, disminuyendo la temperatura del agua en cierta cantidad.
Si le agrego dos hielos, al derretirse, la temperatura final del agua será aún menor.
Lo que por inducción, te hace pensar que debe haber una cantidad óptima de hielo que permite enfriar a su temperatura mínima el agua.
Es decir, que apartir de dicha cantidad de hielo agregado, el agua no se enfriará más.
\section{Dessarrollo}
Definamos primero las cantidades reelevantes para este problema:
\begin{itemize}
  \item Masa de agua $m_a$ en kilogramos
  \item Temperatura inicial del agua $T_a$ en grados centígrados
  \item Masa total del hielo $m_h$ en kilogramos
  \item Temperatura inicial del hielo: $0^\circ C$ (para simplificar)
\end{itemize}
El calor $Q$ que puede \underline{ceder} el agua al enfriarse hasta una temperatura final $T_f$ es:
\begin{equation}
  Q_{agua} = m_a c (T_a - T_f)
  \label{eq:calorEspecifico}
\end{equation}
donde $c \approx 4186\ \mathrm{J/(kg\cdot^\circ C)}$ es el calor específico del agua.

El calor que necesita el hielo para derretirse es 
\begin{equation}
  Q_{fusion} = m_h \cdot L_f
  \label{eq:calorLatente}
\end{equation}
donde $L_f \approx 334000\ \mathrm{J/kg}$ es la constante de fusión del agua.

y si el hielo derretido termina a temperatura $T_f$
\begin{equation}
  Q_{hielo} = m_h L_f + m_h\cdot c \cdot T_f
  \label{eq:calorFinalHielo}
\end{equation}
nótese que el primer término de la suma de la equación \eqref{eq:calorFinalHielo} es la ecuación \eqref{eq:calorLatente} es decir, el calor necesario
para derretir el hielo y el segundo término es la ecuación \eqref{eq:calorEspecifico} es decir, el calor necesario para calentar dicha agua una vez derretida.

El equilibrio térmico dentro del sistema del dispensador de agua con hielos se alcanza cuando
\begin{equation}
  m_a \cdot c \cdot (T_a - T_f) = m_h L_f + m_h \cdot c \cdot T_f
  \label{eq:equilibioTermico}
\end{equation}

Donde el término a la izquierda de la igualdad representa el calor que cede el agua al enfriarse y se transfiere al lado derecho de la igualdad que representa 
el calor que gana el hielo al calentarse.

De la ecuación \eqref{eq:equilibioTermico} puede despejarse la temperatura final del agua $T_f$ para cualquier 
cantidad de agua y masa de hielo.

Por conservación de la energía, no podemos enfriar a menos de $0^\circ C$ agua líquida a presión atmosférica normal usando únicamente hielo.
Por lo que asumiremos que la temperatura mínima alcanzable por el sistema es 
\[ T_f = 0^\circ C \]

\section{Solución}
Ahora, finalmente, podemos plantear y resolver nuestra pregunta principal:
Calcular cuánta masa de hielo se necesita para llevar una masa de agua a exactamente $0^\circ C$

El agua libera 
\begin{equation}
  Q = m_a \cdot c \cdot T_a
  \label{eq:calorAgua}
\end{equation}
Ese calor se usa únicamente para fundir el hielo
\begin{align}
  Q &= m_h \cdot L_f \\
  m_h &= \frac{Q}{L_f} \\
  m_h &= \frac{m_a \cdot c \cdot T_a}{L_f}
  \label{eq:fundicionHielo}
\end{align}
Sabemos que la temperatura del agua potable en Marte, al salir del grifo, es de $20^\circ C$ y que el dispensador
donde se almacena tiene una capacidad de 10 litros.
Sabemos también que 10 litros de agua son equivalentes a 10 kilogramos de agua.
Evaluamos entonces nuestra ecuación:

\begin{equation}
\begin{aligned}
m_h &=
\frac{
(10\,\mathrm{kg})
(4186\,\mathrm{J/(kg\cdot ^\circ C)})
(20\,^\circ\mathrm{C})
}{
334000\,\mathrm{J/kg}
} \\
&\approx 2.5\,\mathrm{kg}
\end{aligned}
\label{eq:calculoMasaHielo}
\end{equation}
Entonces, para enfriar 10 litros de agua a temperatura ambiente en un dispensador con aislamiento térmico a la temperatura mínima
alcanzable, se necesitan 2.5 kilogramos de hielo.

No utilizamos la ley de enfriamiento de Newton ni ninguna otra ecuación diferencial porque para 
este problema sólo nos interesa la temperatura final de equilibrio, no el comportamiento
del sistema en función del tiempo.

\begin{thebibliography}{9}

\bibitem{tippens}
Tippens, Paul E.
\textit{Física: conceptos y aplicaciones}.
Séptima edición revisada.
McGraw-Hill, 2001.
pp.~354--355.

\end{thebibliography}
\end{document}
